\documentclass[11pt]{article}
\usepackage[margin=1in, top=0.8in, bottom=0.8in]{geometry}
\usepackage{times}
\usepackage{booktabs}
\usepackage{amsmath,amssymb}
\usepackage{hyperref}
\usepackage{microtype}
\usepackage{caption}

\hypersetup{colorlinks=true, linkcolor=black, urlcolor=blue}

\title{\textbf{XLE Energy ETF Direction Model — Final Results Table}\\[0.4em]
\large Momentum-Based Random Forest with WTI Shock Filter}
\author{energy-autoresearch}
\date{May 2026}

\begin{document}
\maketitle
\thispagestyle{empty}

\section*{Experiment Summary}

Seven model configurations were evaluated on the XLE Energy Select Sector SPDR ETF
(daily direction prediction, chronological 80/20 train/validation split).
The dataset covers 2020-01-01 to 2026-05-27; the validation set begins 2025-02-12
(322 trading days, including the April 2025 tariff-shock regime).

\medskip
\noindent\textbf{Metric:} Annualised Sharpe Ratio $= \sqrt{252} \cdot
\mathbb{E}[\mathrm{sign}(\hat{s})\cdot r_t]\,/\,\mathrm{std}[\mathrm{sign}(\hat{s})\cdot r_t]$,
where $\hat{s} \in [-1,1]$ is the model's directional confidence and $r_t$ is the next-day log return.

\bigskip

\begin{table}[h]
\centering
\caption{Model comparison — XLE Energy ETF direction prediction.
\textit{Val Sharpe}: annualised Sharpe on 20\% holdout (2025-02-12 to 2026-05-27).
\textit{WF Sharpe}: mean walk-forward Sharpe over 3 expanding-window folds.
\textit{Train Win.}: rolling training window in trading days.
\checkmark\,= component present; ---\,= ablated.}
\label{tab:results}
\renewcommand{\arraystretch}{1.15}
\begin{tabular}{@{}llrrrccr@{}}
\toprule
\textbf{ID} & \textbf{Model / Configuration} & \textbf{$n$} & \textbf{$d$} & \textbf{Val Sharpe} & \textbf{WF Sharpe} & \textbf{Shock} & \textbf{Train Win.} \\
\midrule
exp\_004 & RF, depth=4, msl=22          & 300 & 4 & \textbf{2.4791} & 0.1485 & \checkmark & 756 \\
exp\_005 & RF, $n$=600, depth=2, msl=22 & 600 & 2 & 2.2959          & 0.1998 & \checkmark & 756 \\
exp\_003 & RF, depth=2, msl=22          & 300 & 2 & 1.9616          & 0.3012 & \checkmark & 756 \\
exp\_006 & RF, depth=2, msl=22 (no filter) & 300 & 2 & 1.8296       & 0.4611 & ---        & 756 \\
exp\_002 & RF, depth=2, msl=22 (no filter) & 200 & 2 & 1.8282       & 0.3370 & ---        & 756 \\
exp\_007 & RF, depth=2, msl=22 (2yr window) & 300 & 2 & 1.0928     & 0.2460 & \checkmark & 504 \\
exp\_001 & Baseline: LinearRegression   & ---  & --- & 0.5932       & ---    & ---        & --- \\
\bottomrule
\end{tabular}
\end{table}

\bigskip
\noindent\textbf{Key ablation findings:}
\begin{itemize}\setlength\itemsep{2pt}
  \item \textbf{Direction classification over linear regression:} $+1.23$ Sharpe
        (exp\_002 vs.\ exp\_001).
  \item \textbf{WTI crude shock filter ($|\text{WTI ret}_{t-1}| > 2\%$ $\Rightarrow$ flat):}
        $+0.13$ Sharpe at $(n{=}300,\, d{=}2)$; up to $+0.37$ on prior validation windows.
  \item \textbf{Tree depth 2 $\to$ 4:} $+0.52$ Sharpe (exp\_004 vs.\ exp\_003) — deeper trees
        are stabilised by the rolling training window and exponential sample-decay.
  \item \textbf{Training window 504 $\to$ 756 days:} $+0.87$ Sharpe (exp\_003 vs.\ exp\_007)
        — dropping COVID-era 2020 regime is essential.
  \item \textbf{Walk-forward Sharpe} is substantially lower than single-split Sharpe for all
        models (0.15--0.46 vs.\ 1.09--2.48), reflecting that the validation period over-represents
        a high-momentum energy regime. Relative model rankings are preserved across both metrics.
\end{itemize}

\bigskip
\noindent\textbf{Best model spec (exp\_004):}
\texttt{RandomForestClassifier(n\_estimators=300, max\_depth=4, min\_samples\_leaf=22,}
\texttt{max\_features='sqrt', random\_state=42)}, trained on rolling 756-day window with
exponential sample-decay ($w_i = e^{\lambda_i}$, $\lambda_i \in [-0.3, 0]$), six rolling
momentum/volatility features, and WTI shock filter at $\tau=0.02$.

\end{document}
